% Part: normal-modal-logic
% Chapter: syntax-and-semantics
% Section: normal-modal-logics

\documentclass[../../../include/open-logic-section]{subfiles}

\begin{document}

\olfileid{nml}{syn}{nor}

\olsection{సాధారణ మోడల్ తర్కాలు}

మోడల్ తర్కాన్ని తగిన ధర్మాలు కలిగిన \tetoken{సూత్రాల}{!!{formula}s} ఏ సమితిగానైనా చూడవచ్చు. ఆసక్తికరమైన అన్వయాలు, ఉపయోగాలు ఉండటం వల్లే కాక, ముఖ్యంగా సిద్ధాంతపరంగా బాగా అర్థమైనవి కావడం వల్ల కొన్ని మోడల్ తర్కాలకు ప్రత్యేక ప్రాముఖ్యం ఉంది. వాటినే సాధారణ మోడల్ తర్కాలంటారు.

\begin{defn}\ollabel{defn:modal-logic}
ప్రాథమిక మోడల్ భాషలోని \tetoken{సూత్రాల}{!!{formula}s} సమితి~$\Log L$ను \emph{మోడల్ తర్కం} అంటాం, ఒకవేళ అది
\begin{enumerate}
\item అన్ని ప్రతిపాదనాత్మక సర్వసత్యాలను కలిగి ఉంటే,
\item ఏకరీతి ప్రతిస్థాపన కింద సంవృతమైతే, మరియు
\item మోడస్ పోనెన్స్ కింద సంవృతమైతే; అంటే $!A$, $(!A \lif
  !B) \in \Log L$ అయినప్పుడల్లా $!B \in \Log L$ కూడా అయితే.
\end{enumerate}
\end{defn}

\begin{ex}
ఈ నిర్వచనాన్ని తీర్చేవి, అందువల్ల మోడల్ తర్కాలుగా పరిగణించబడేవి, అంత ఆసక్తికరంకాని కొన్ని \tetoken{సూత్రాల}{!!{formula}s} సమితులు:
\begin{enumerate}
\item అన్ని సర్వసత్యాల సమితి,
\item అన్ని \tetoken{సూత్రాల}{!!{formula}s} సమితి (వైరుధ్యపూరిత మోడల్ తర్కం).
\end{enumerate}
\end{ex}

\begin{defn}\label{defn:normal-modal-logic}
మోడల్ తర్కం~$\Log L$ \emph{సాధారణమైనది} కావడానికి, కావలసినవీ చాలునవీ:
\begin{enumerate}
\item అది ఈ రెండింటినీ కలిగి ఉండాలి:
\begin{align}
\tag{K} \Box(p \lif q) & \lif (\Box p \lif \Box q) \\
\tag{Dual} \Diamond p & \liff \lnot \Box \lnot p
\end{align}
\item అలాగే అనివార్యతీకరణ కింద సంవృతమై ఉండాలి; అంటే $!A \in
  \Log L$ అయినప్పుడల్లా $\Box !A \in \Log L$ కూడా కావాలి.
\end{enumerate}
\end{defn}

మోడల్ తర్కాలను నిర్వచించే రెండు ప్రధాన మార్గాలను చూస్తాం. ఒకటి నిరూపణ-సిద్ధాంతపరంగా: నిర్దిష్ట \tetoken{వ్యుత్పత్తి}{!!{derivation}} వ్యవస్థల్లో వ్యుత్పాదించగల \tetoken{సూత్రాల}{!!{formula}s} సమితిగా. మరొకటి నమూనా-సిద్ధాంతపరంగా: నిర్దిష్ట నమూనాల వర్గాల్లో చెల్లుబాటయ్యే \tetoken{సూత్రాల}{!!{formula}s} సమితులుగా. ఇది సాధారణ ప్రతిపాదనాత్మక తర్కంలో (లేదా మొదటిస్థాయి తర్కంలో) పరిస్థితిని ప్రతిబింబిస్తుంది. అక్కడ కూడా \tetoken{సూత్రాల}{!!{formula}s} సమితిని రెండు విధాలుగా వేరు చేయవచ్చు: అన్ని సత్యమూల్య కేటాయింపుల్లో సత్యమయ్యే సర్వసత్యాల సమితిగా, లేదా ఏదో \tetoken{వ్యుత్పత్తి}{!!{derivation}} వ్యవస్థలోని అన్ని \tetoken{వ్యుత్పాదించదగిన}{!!{derivable}} \tetoken{సూత్రాల}{!!{formula}s} సమితిగా. సాధారణ మోడల్ తర్కాల విషయంలో ఈ ద్వంద్వ నిర్దేశం సంబంధ నమూనాలు, స్వీకృత-ఆధారిత (హిల్బర్ట్ తరహా) నిరూపణ వ్యవస్థల ద్వారా సాధ్యం.
\end{document}
